\section{机器人感知}

\subsection{感知的独特挑战}

机器人感知面临的挑战远超传统计算机视觉：

\begin{figure}[H]
\centering
\includegraphics[width=0.95\textwidth]{slides-images/slide-009.jpg}
\caption{机器人感知面临的挑战：不完整信息、遮挡、传感器误差、动态环境、可形变物体等\protect\footnotemark}
\end{figure}
\footnotetext{Slides 第10页。}

\begin{itemize}
    \item 观测只包含环境的\textbf{不完整信息}（遮挡、部分可观测）
    \item 传感器数据存在噪声和误差
    \item 动作执行可能失败（如抓取滑落）
    \item 环境包含刚体、可形变物体（布料、液体）、甚至其他智能体（人、动物）
    \item 环境是动态变化的
\end{itemize}

\subsection{多模态感知}

\begin{figure}[H]
\centering
\includegraphics[width=0.95\textwidth]{slides-images/slide-010.jpg}
\caption{机器人的多模态传感器：视觉（RGB/RGBD）、触觉、听觉、深度传感器协同工作\protect\footnotemark}
\end{figure}
\footnotetext{Slides 第11页。}

在机器人领域，仅依赖摄像头远远不够。研究者会尽可能多地为机器人配备传感器：
\begin{itemize}
    \item \textbf{触觉传感}：判断抓取是否稳定
    \item \textbf{音频信息}：感知物理接触事件
    \item \textbf{深度信息}：获取三维几何
    \item \textbf{摄像头}：宏观环境状态
\end{itemize}

这些传感器互补协作，为机器人提供全面的环境理解。

\subsection{机器人视觉的三大特征}

\begin{figure}[H]
\centering
\includegraphics[width=0.95\textwidth]{slides-images/slide-012.jpg}
\caption{机器人视觉的三大特征：具身性（Embodied）、主动性（Active）、情境性（Situated）\protect\footnotemark}
\end{figure}
\footnotetext{Slides 第13页。}

\begin{importantbox}{机器人视觉 $\neq$ 计算机视觉}
\begin{enumerate}
    \item \textbf{具身性（Embodied）}：机器人有物理躯体，直接在物理世界中体验，动作会即时影响自身的感知
    \item \textbf{主动性（Active）}：机器人是主动感知者，它知道自己想感知什么，能选择何时、何处、如何感知（如移动头部看桌子后方）
    \item \textbf{情境性（Situated）}：机器人处于真实世界中，感知必须与当前任务紧密耦合——扣扣子时只需关注按钮局部区域
\end{enumerate}
\end{importantbox}

\begin{figure}[H]
\centering
\includegraphics[width=0.95\textwidth]{slides-images/slide-011.jpg}
\caption{计算机视觉中的实例分割（左）vs 机器人视觉中的主动交互感知（右）：机器人通过推动物体来判断是一个整体还是多个堆叠部件\protect\footnotemark}
\end{figure}
\footnotetext{Slides 第12页。}

\subsection{本章小结}

机器人感知不同于被动的计算机视觉。它是具身的、主动的、情境化的。机器人需要多模态传感器协同工作，并且感知系统需要与下游决策系统紧密耦合，实现闭环的感知-行动循环。

%% ========== Part 3: 强化学习 ==========

